% chapters/04/section-03.tex
\section{Mengendalikan Laju dengan \istilah{Clock}}
\label{sec:mengendalikan-laju}

Siklus permainan yang berjalan tanpa pengendali waktu akan memakai seluruh
kemampuan prosesor untuk menghasilkan bingkai sebanyak mungkin. Cara itu
merugikan, dan bagian ini menjelaskan perangkat yang mencegahnya.

\subsection{Dua Pekerjaan \code{pygame.time.Clock}}

Pygame menyediakan perkakas pengendali waktu bernama \code{Clock}. Satu
pemanggilan perintahnya, yaitu \code{clock.tick()}, mengerjakan dua hal
sekaligus.

\begin{enumerate}
  \item \textbf{Mengukur selang waktu} yang benar-benar berlangsung sejak
        pemanggilan sebelumnya, dan mengembalikan nilainya dalam milidetik.
        Nilai inilah yang pada Bab~V dipakai untuk menghitung gerak.
  \item \textbf{Menahan laju} dengan menunggu sampai selang waktu satu bingkai
        terpenuhi, apabila batas laju diberikan sebagai argumen. Apabila
        seluruh pekerjaan satu bingkai selesai lebih cepat daripada jatah
        waktunya, \code{tick()} menunda sisa waktunya.
\end{enumerate}

Karena kedua pekerjaan itu digabung dalam satu perintah, \code{tick()} selalu
diletakkan pada akhir setiap putaran siklus, yaitu setelah tahap penayangan.
Peletakan itu penting: selang waktu yang diukur harus mencakup seluruh
pekerjaan satu bingkai, bukan sebagiannya.

\begin{kodekeluaran}
delta_seconds = clock.tick(60) / 1000.0
\end{kodekeluaran}

Batas laju dinyatakan dalam bingkai per detik. Nilai 60 berarti paling banyak
60 bingkai setiap detik. Perhatikan pembagian dengan 1000: \code{tick()}
mengembalikan milidetik, sedangkan seluruh perhitungan gerak nanti memakai
satuan detik. Kebiasaan mengubahnya menjadi detik sejak awal mencegah
kekeliruan satuan yang sulit dilacak.

\begin{catatabox}
Batas laju bersifat pilihan, bukan keharusan. Memanggil
\code{clock.tick()} tanpa argumen tidak menahan laju sama sekali, sedangkan
\code{clock.tick(0)} berarti hal yang sama. Cara itu hanya berguna untuk
mengukur seberapa cepat sebenarnya komputer Anda menyelesaikan satu bingkai.
\end{catatabox}

\subsection{Mengukur Laju yang Benar-Benar Tercapai}

Batas laju menyatakan keinginan, bukan kenyataan. Apabila pekerjaan satu
bingkai memerlukan waktu lebih lama daripada jatahnya, laju yang tercapai
akan lebih rendah daripada batas yang diminta. Karena itu Pygame menyediakan
cara mengukur laju yang benar-benar terjadi.

\begin{kodekeluaran}
measured_fps = clock.get_fps()
\end{kodekeluaran}

Nilai \code{clock.get_fps()} dihitung dari rata-rata selang waktu beberapa
bingkai terakhir, sehingga angkanya tidak melompat-lompat. Nilai itu berguna
untuk dua hal: menampilkan keadaan program kepada pemain, dan menolong Anda
sebagai penulis program. Perbedaan yang besar antara batas laju dan laju
terukur merupakan tanda paling awal bahwa program Anda terlalu berat.

Pada program bab ini, nilai tersebut ditampilkan pada judul jendela setiap
bingkai. Cara itu dipilih karena judul jendela selalu terlihat, tidak
memakan ruang gambar permainan, dan dapat dibaca walaupun jendela sedang
tidak difokuskan.

\subsection{Mengapa Laju Bingkai Perlu Dibatasi}

Pertanyaan yang wajar muncul: bukankah permainan yang lebih cepat lebih baik?
Jawabannya tidak, dan alasannya ada tiga.

\begin{enumerate}
  \item \textbf{Perilaku permainan menjadi tidak seragam.} Apabila gerak
        bergantung pada jumlah bingkai, permainan berjalan lebih sulit pada
        komputer yang lebih cepat. Batas laju membuat seluruh komputer
        menghasilkan jumlah bingkai yang sama setiap detik, sehingga
        permainan terasa sama di mana pun.
  \item \textbf{Beban prosesor dan panas berlebihan.} Prosesor yang dipakai
        penuh tanpa manfaat menaikkan suhu, mempercepat pemakaian baterai pada
        komputer jinjing, dan memperpendek usia perangkat. Angka pengukuran
        pada Bagian~\ref{sec:mengukur-laju} memperlihatkan besarnya
        perbedaan itu.
  \item \textbf{Tidak ada manfaat yang sepadan.} Mata manusia tidak dapat
        membedakan dua ribu bingkai per detik, dan layar hanya dapat
        menampilkan sejumlah gambar terbatas setiap detik. Bingkai di atas
        batas kemampuan layar terbuang percuma.
\end{enumerate}

Sebagai gambaran awal, siklus yang dibatasi pada 60 bingkai per detik memakai
sekitar tiga persen dari satu inti prosesor, sedangkan siklus yang sama tanpa
pembatasan memakai hampir seratus persen untuk pekerjaan yang tidak terlihat
pemain. Perbandingan lengkapnya beserta cara mengukurnya sendiri dibahas pada
Bagian~\ref{sec:mengukur-laju}.

\begin{tipsbox}
Nilai batas laju yang wajar untuk pembelajaran adalah 60 bingkai per detik,
karena nilai itu mulus bagi mata dan cukup ringan bagi sebagian besar
komputer. Turunkan menjadi 30 apabila program Anda mulai berat, dan
naikkan hanya apabila Anda mengetahui alasan yang jelas untuk itu.
\end{tipsbox}
